% Tableau genere automatiquement depuis bench/results/human_e2e.json - ne pas editer a la main
\section{Annexe : corpus de la contre-épreuve à locuteur humain}
\label{sec:annexe-humain}

Le tableau~\ref{tab:annexe-humain} donne le corpus intégral de la contre-épreuve à
locuteur humain réel (section~\ref{sec:res-e2e}) : les 25 clairances dans leur forme
écrite du corpus annoté, l'ordre TrafScript attendu après validation, et le verdict
d'exécution exacte dans les deux conditions (prises micro brutes ; canal radio simulé,
bruit SNR 12~dB puis passe-bande VHF). Les clairances écrites en abrégé
(\texttt{AFR1234}, \texttt{FL320}, \texttt{hdg}) ont été prononcées en phraséologie
parlée (\emph{air france one two three four}, \emph{flight level three two zero},
\emph{heading}). Ce tableau est généré automatiquement depuis les données brutes
versionnées (\texttt{bench/results/human\_e2e.json}), comme toutes les valeurs
chiffrées du document.

\begin{table}[h]
\centering
\footnotesize
\begin{tabular}{>{\raggedright}p{0.5cm} >{\raggedright}p{6.6cm} >{\raggedright}p{4.4cm} c c}
\toprule
\rowcolor{lightgray}\textbf{\#} & \textbf{Clairance (forme écrite du corpus)} & \textbf{Ordre attendu (TrafScript)} & \textbf{Micro} & \textbf{Radio}\\
\midrule
1 & air france one two three four turn left heading two seven zero & \texttt{HDG AFR1234 270} & OK & OK\\
2 & speedbird five seven turn right heading 090 & \texttt{HDG BAW57 90} & OK & OK\\
3 & AFR1234 fly heading three six zero & \texttt{HDG AFR1234 360} & OK & KO\\
4 & EZY21 fly heading 100 & \texttt{HDG EZY21 100} & OK & OK\\
5 & lufthansa eight eight heading one eight zero & \texttt{HDG DLH88 180} & OK & OK\\
6 & ryanair niner turn right heading two two five & \texttt{HDG RYR9 225} & KO & KO\\
7 & AFR1234 hdg 045 & \texttt{HDG AFR1234 45} & OK & OK\\
8 & BAW57 turn left heading 310 & \texttt{HDG BAW57 310} & OK & OK\\
9 & air france one two three four descend flight level one zero zero & \texttt{ALT AFR1234 10000} & OK & OK\\
10 & csa one delta zulu climb flight level two four zero & \texttt{ALT CSA1DZ 24000} & KO & KO\\
11 & AFR1234 climb FL320 & \texttt{ALT AFR1234 32000} & OK & OK\\
12 & BAW57 descend to flight level 80 & \texttt{ALT BAW57 8000} & OK & OK\\
13 & EZY21 climb 5 thousand feet & \texttt{ALT EZY21 5000} & OK & KO\\
14 & AFR1234 descend altitude 6000 feet & \texttt{ALT AFR1234 6000} & OK & OK\\
15 & air france one two three four climb to flight level three two zero & \texttt{ALT AFR1234 32000} & OK & OK\\
16 & BAW57 maintain flight level 100 & \texttt{ALT BAW57 10000} & OK & OK\\
17 & DLH88 descend 4000 feet & \texttt{ALT DLH88 4000} & OK & OK\\
18 & speedbird five seven climb flight level one five zero & \texttt{ALT BAW57 15000} & OK & OK\\
19 & AFR1234 reduce speed two five zero knots & \texttt{SPD AFR1234 250} & KO & KO\\
20 & AFR1234 increase speed 300 & \texttt{SPD AFR1234 300} & OK & OK\\
21 & RYR9 increase speed to 280 knots & \texttt{SPD RYR9 280} & KO & KO\\
22 & easyjet two one reduce speed two two zero & \texttt{SPD EZY21 220} & OK & OK\\
23 & BAW57 speed 250 & \texttt{SPD BAW57 250} & OK & OK\\
24 & csa one delta zulu reduce 230 & \texttt{SPD CSA1DZ 230} & KO & KO\\
25 & AFR1234 maintain 250 knots & \texttt{SPD AFR1234 250} & OK & OK\\
\midrule
 & \multicolumn{2}{l}{\textbf{Total d'exécutions exactes}} & \textbf{20/25} & \textbf{18/25}\\
\bottomrule
\end{tabular}
\caption{Contre-épreuve à locuteur humain réel : verdict phrase par phrase dans les
deux conditions. OK = ordre exécuté exactement conforme à l'attendu ; KO = écart
(transcription déviée ou interprétation erronée, détail en
section~\ref{sec:res-e2e}).}
\label{tab:annexe-humain}
\end{table}
